\section{Conclusion}

We presented Delta, a domain-specific language whose declarative primitives reify the recurring patterns we identified across 46 existing interactive formula instances into a coherent authoring model. In two user studies, Delta made formula augmentation faster and easier than a conventional React approach, and in open-ended tasks its declarative structure scaffolded not only implementation but also pedagogical reasoning about what to make interactive and why. More broadly, Delta demonstrates that a well-chosen set of domain-specific abstractions can collapse the distance between mathematical intent and interactive artifact — a principle likely to generalize to other domains where authors want to make formal notation explorable. We hope this work opens new lines of inquiry into DSL-driven authoring tools, graphical interfaces built atop declarative specifications, and intelligent support that shifts the bottleneck from implementation to pedagogical design.